%%XELATEX- UTF8%%
\documentclass[french]{book}

\usepackage[T1]{fontenc}
\usepackage[utf8]{inputenc}
\usepackage[french]{babel}
\usepackage{teach}
\usepackage{verbatim}
\usepackage{amsmath,amsfonts,amssymb}
\usepackage{pgfplots}
\RequirePackage{graphicx}
\RequirePackage{ifpdf}
 \ifpdf
   \DeclareGraphicsRule{*}{mps}{*}{}
 \fi
\usepackage{fancyhdr}
\usepackage{pstricks}
\usepackage{amsmath}
\usepackage{amssymb}
\usepackage{amsfonts}
\usepackage{amsthm}
\usepackage{mathrsfs}
\usepackage{array}
\usepackage{multicol}
\setlength{\multicolsep}{3pt}
\usepackage{graphicx}
\graphicspath{{Figures/}}
\usepackage{pstricks,pst-plot,pst-text,pst-tree,pst-eucl}
\usepackage[all]{xy}
\usepackage{fancybox}
\usepackage{color}
\usepackage{hyperref}
%  \usepackage[frenchb]{babel}
\usepackage{subfig}
\pgfplotsset{compat=newest}
%\usepackage{geometry}
%\geometry{top=1cm, bottom=1cm, left=2cm, right=2cm}
%\usepackage[utf8]{inputenc}
\usepackage[french]{babel}
\redefinecolor{thm}{title}{bgcolor}{alizarin}
\redefinecolor{prop}{title}{bgcolor}{meatbrown}
\redefinecolor{defn}{title}{bgcolor}{olivine}
\definecolor{alizarin}{rgb}{0.82, 0.1, 0.26}
\definecolor{meatbrown}{rgb}{0.9, 0.72, 0.23}
\definecolor{olivine}{rgb}{0.6, 0.73, 0.45}
\usepackage{mathrsfs}
%%
\usepackage{yhmath} 
\newcommand{\R}{\mathbb {R}}
%\newcommand{\C}{\mathds {C}}
\newcommand{\N}{\mathbb {N}}
\newcommand{\Z}{\mathbb {Z}}
\newcommand{\Q}{\mathbb {Q}}
\newcommand{\D}{\mathbb {D}}
%%
\newcommand{\se}{\geq}
\newcommand{\ie}{\leq}
\newcommand{\tm}{\times}

%\newcommand{\ell}{l}
%\newcommand{\ell'}{l'}
%%

\newcommand{\sysd}[2]{%
       \left\{
       \begin{aligned}
          #1\\
          #2\\
       \end{aligned}
       \right.%
       }

\newcommand{\ds}{\displaystyle}
\newcommand{\limn}{\lim_{n\rightarrow +\infty}}
\newcommand{\limo}[1][x]{\ds\lim_{#1\rightarrow 0}}
\newcommand{\limite}[2][x]{\ds\lim_{#1\rightarrow #2}}
\newcommand{\limpinf}[1][x]{\ds\lim_{#1\rightarrow +\infty}}
\newcommand{\limminf}[1][x]{\ds\lim_{#1\rightarrow -\infty}}
\newcommand{\limiteg}[2][x]{\ds\lim_{#1\xrightarrow{<}#2}}
\newcommand{\limited}[2][x]{\ds\lim_{#1\xrightarrow{>}#2}}
%%
\newcommand{\vect}[1]{\overrightarrow{#1}}
\newcommand{\arc}[1]{\wideparen{#1}}
\newcommand{\oij}{(\textit{O},{\overrightarrow{i}},{\overrightarrow{j}})}
%\newcommand{\imath}{\overrightarrow{i}}
%\newcommand{\jmath}{\overrightarrow{j}}
%
\newcommand{\cp}[2]{%
        \begin{pmatrix}
        #1\\
        #2
        \end{pmatrix}%
        }
%
\newcommand{\calb}{\mathscr{B}}
\newcommand{\calc}{\mathscr{C}}
\newcommand{\cald}{\mathscr{D}}
\newcommand{\cale}{\mathscr{E}}
\newcommand{\calf}{\mathscr{F}}
\newcommand{\calp}{\mathscr{P}}
\newcommand{\cals}{\mathscr{S}}
\newcommand{\calt}{\mathscr{T}}
%
\newcommand{\suite}[1][u]{\left( #1_{n} \right)_{n\in\N}}
\newcommand{\suitep}[1][u]{\left( #1_{n} \right)_{n\in\N^{*}}}
\usetikzlibrary{arrows}

\newcommand{\poly}[1][a]{#1_0+#1_1x+\cdots+#1_nx^n}

\begin{document}
\tableofcontents
\setcounter{chapter}{4}
\chapter{Trigonom\'etrie}

\section{Angles orientés}
Une unité de longueur est choisie.
\subsection{Cercle trigonométrique, mesures d'un arc}
\begin{minipage}{10cm}
\begin{defn}
On appelle cercle trigonométrique un cercle de rayon 1, muni d'une origine et orienté dans le sens inverse des aiguilles d'une montre. Ce sens est appelé sens direct ou sens trigonométrique.
\par Soit $\calc$ un cercle trigonométrique de centre $O$ et d'origine $A$. é tout réel $x$, on peut associer un unique point $M$ du cercle en \og enroulant\fg{} la droite des réels autour du cercle. On dit que $x$ est une mesure de l'arc orienté $\arc{AM}$.
\end{defn}

\begin{rem}
Un point correspond à une infinité de réels.
\end{rem}

\begin{exemple}
Le point $A$ est associé aux réels 0, $2\pi$, $4\pi$, $-2\pi$...
Le point $B$ est associé aux réels $\dfrac{\pi}{2}$, $\dfrac{5\pi}{2}$, $-\dfrac{3\pi}{2}$...
Le point $A'$ est associé aux réels $\pi$, $3\pi$, $-\pi$...
\end{exemple}
\end{minipage}
\hfill
\begin{minipage}{4.6cm}
\psset{unit=1.3cm,PointSymbol=none}
\begin{center}
\begin{pspicture}(-2,-3)(1.5,3)
\pstGeonode[PosAngle=-135]{O}
\pstGeonode[PosAngle={0,90,180,-90},PointNameSep=0.5em](1,0){A}(0,1){B}(-1,0){A'}(0,-1){B'}
\pscircle(0,0)1
\pstLineAB{A'}{A}
cc\pstLineAB{B'}{B}
\pstGeonode[PointName=none](0.707,0.707){a}(-1.2,2.614){b}(-0.707,0.707){c}(-2,-0.586){d}
\pstOrtSym[PointName=none]{A'}{A}{a,b,c,d}
\psline{->}(1,-3)(1,3)
%\psline[linestyle=dashed]{->}(1,-3)(1,3)
\pstLineAB[linestyle=dashed]{->}{a}{b}
\pstLineAB[linewidth=1.8\pslinewidth]{->}{c}{d}
\pstLineAB[linestyle=dashed]{a'}{b'}
\pstLineAB[linewidth=1.8\pslinewidth]{c'}{d'}
\pstArcOAB[linewidth=1.8\pslinewidth]{O}{c'}{c}
\pstGeonode[PointSymbol=|,dotangle=90,dotscale=1.3](1,2.094){x}
\pstGeonode[PointSymbol=|,dotangle=90,dotscale=1.3,PointName=1](1,1){i}
\pstGeonode[PointSymbol=|,dotangle=-45,dotscale=1.3,PointName=x,PosAngle=15](-0.218,1.633){x'}
\pstGeonode[PointSymbol=|,dotangle=30,dotscale=1.3,PosAngle=150](-0.5,0.866){M}
\ncarc[nodesep=3mm,arcangle=-30]{->}{x}{x'}
\ncarc[nodesep=2mm,arcangle=-30]{->}{x'}{M}
\end{pspicture}
\end{center}
\end{minipage}

\subsection{Mesures d'un angle orienté de vecteurs}
étant donnés deux vecteurs non nuls $\vect{u}$ et $\vect{v}$ et un cercle trigonométrique $\calc$ de centre $O$, on appelle $A$ le point d'intersection de $\calc$ et de la droite $(O,\vect{u})$ et $B$ le point d'intersection de $\calc$ et de la droite $(O,\vect{v})$.

\begin{minipage}{10.5cm}
\begin{defn}
On appelle mesure de l'angle orienté $(\vect{u},\vect{v})$ toute mesure de l'arc orienté $\arc{AB}$.

Si $x$ est une de ces mesures, toute autre mesure s'écrit $y=x+2k\pi$ avec $k\in\Z$.

On note (de façon abusive) : $(\vect{u},\vect{v})=x+k\times 2\pi, k\in\Z$.
\end{defn}
\end{minipage}
\hfill
\begin{minipage}{4cm}
	\begin{center}
	\includegraphics{CoursAnglesFigures.1}
	\end{center}
\end{minipage}

% Remarque : Si $x$ est une mesure de $(\vect{u},\vect{v})$, toute autre mesure s'�crit $y=x+2k\pi$ avec $k\in\Z$. On note donc aussi $(\vect{u},\vect{v})=x+2k\pi, \; k\in\Z$.

\begin{defn}
On appelle mesure principale de $(\vect{u},\vect{v})$, l'unique mesure appartenant à $]-\pi ; \pi]$.

\vspace{4mm}
\end{defn}

\begin{exemple}
Déterminer la mesure principale d'un angle dont une mesure est $\dfrac{126\pi}{5}$
\end{exemple}


\[\dfrac{126\pi}{5} = \dfrac{12\times10\pi+6\pi}{5} = 12\times2\pi+\dfrac{6\pi}{5} = 12\times2\pi+2\pi-\dfrac{4\pi}{5} = -\dfrac{4\pi}{5}+13\times 2\pi\]
La mesure principale de cet angle est donc $\dfrac{-4\pi}{5}$.


\section{Propriétés}
\subsection{Angles et colinéarité}
\begin{prop}
Soient $\vect{u}$ et $\vect{v}$ deux vecteurs non nuls.
\par $\vect{u}$ et $\vect{v}$ sont colinéaires et de même sens si et seulement si $(\vect{u},\vect{v})=0+2k\pi$.
\par $\vect{u}$ et $\vect{v}$ sont colinéaires et de sens contraires si et seulement si $(\vect{u},\vect{v})=\pi+2k\pi$.
\end{prop}

\begin{center}
\includegraphics{CoursAnglesFigures.2}
\hspace{3em}
\includegraphics{CoursAnglesFigures.3}
\end{center}


\subsection{Relation de Chasles}
\begin{minipage}{10cm}
\begin{prop}
Pour tous vecteurs non nuls $\vect{u}$, $\vect{v}$ et $\vect{w}$ : \[(\vect{u},\vect{v})+(\vect{v},\vect{w})=(\vect{u},\vect{w})+2k\pi\]
\end{prop}
\end{minipage}
\hfill
\begin{minipage}{4cm}
\includegraphics{CoursAnglesFigures.4}
\end{minipage}


\begin{exemple}
Dans la figure suivante, démontrer que les droites $(AB)$ et $(CD)$ sont parallèles.
\end{exemple}

\begin{minipage}{9.5cm}

\begin{multline*}
(\vect{AB},\vect{CD})=(\vect{AB},\vect{BC})+(\vect{BC},\vect{CD})+2k\pi\\
=\dfrac{\pi}{3}+\dfrac{2\pi}{3}+2k\pi=\pi+2k\pi
\end{multline*}
\par Les vecteurs $\vect{AB}$ et $\vect{CD}$ sont donc colinéaires.
\par \underline{Conclusion} : $(AB)$ et $(CD)$ sont parallèles.

\end{minipage}
\hfill
\begin{minipage}{5cm}
\begin{center}
\includegraphics{CoursAnglesFigures.9}
\end{center}
\end{minipage}

\newpage

\begin{prop}
\par Pour tous vecteurs non nuls $\vect{u}$ et $\vect{v}$ :

\vspace{1ex}
\begin{tabular}{ll@{\hspace{5em}}ll}
$(\vect{v},\vect{u})=-(\vect{u},\vect{v})+2k\pi$ & $(1)$ & $(\vect{u},-\vect{v})=(\vect{u},\vect{v})+\pi+2k\pi$ & $(2)$ \\[1ex]
$(-\vect{u},\vect{v})=(\vect{u},\vect{v})+\pi+2k\pi$ & $(3)$ & $(-\vect{u},-\vect{v})=(\vect{u},\vect{v})+2k\pi$ & $(4)$
\end{tabular}
\end{prop}

\includegraphics{CoursAnglesFigures.5}
\hfill
\includegraphics{CoursAnglesFigures.6}
\hfill
\includegraphics{CoursAnglesFigures.7}
\hfill
\includegraphics{CoursAnglesFigures.8}

\begin{demo}
\par (1) : $(\vect{v},\vect{u})+(\vect{u},\vect{v})=(\vect{v},\vect{v})+2k\pi=2k\pi$ ainsi $(\vect{v},\vect{u})=-(\vect{u},\vect{v})+2k\pi$
\par (2) : $(\vect{u},-\vect{v})=(\vect{u},\vect{v})+(\vect{v},-\vect{v})+2k\pi=(\vect{u},\vect{v})+\pi+2k\pi$
\par (3) se démontre comme (2).

\begin{tabbing}
(4) : $(-\vect{u},-\vect{v})$ \= $=(-\vect{u},\vect{v})+\pi+2k\pi$ d'après (2)\\
 \> $= (\vect{u},\vect{v})+\pi+\pi+2k\pi$ d'après (3)\\
ainsi $(-\vect{u},-\vect{v})=(\vect{u},\vect{v})+2k\pi$
\end{tabbing}
\vspace{-1em}
\end{demo}


\section{Lignes trigonom\'etriques}

\subsection{Cosinus et Sinus}
\subsubsection*{Rappels}
\begin{minipage}{10cm}
Soit $\oij$ un rep\`ere orthonormal direct (c'est-à-dire tel que $(\vect{\imath},\vect{\jmath})=\frac{\pi}{2}+2k\pi$), $\calc$ le cercle trigonom\'etrique de centre $O$.
\par $A$ et $B$ sont les points tels que $\vect{OA}=\vect{\imath}$ et $\vect{OB}=\vect{\jmath}$.

\vspace{1ex}
\begin{defn}
Pour tout r\'eel $x$, il existe un unique point $M$ de $\calc$ tel que $x$ soit une mesure de l'angle $(\vect{\imath},\vect{OM})$.
\par $\cos x$ est l'abscisse de $M$ dans $\oij$.
\par $\sin x$ est l'ordonn\'ee de $M$ dans $\oij$.
\end{defn}
\end{minipage}
\hfill
\begin{minipage}{4cm}
\includegraphics[scale=1.3]{CoursTrigoFigures.1}
\end{minipage}

\subsubsection*{Propri\'et\'es imm\'ediates}
\begin{prop}
Pour tout $x\in\R$ :

\begin{tabular}{*{3}{m{0.3\linewidth}}}
$-1\ie \cos x \ie 1$ \par $-1\ie \sin x \ie 1$ & $\cos(x+2k\pi)=\cos x$ \par $\sin(x+2k\pi)=\sin x$ & $\cos^2 x+\sin^2 x=1$
\end{tabular}
\end{prop}

\subsubsection*{Cosinus et sinus d'un angle orient\'e de vecteurs}
Soit $(\vect{u},\vect{v})$ un angle orient\'e et $x$ une de ses mesures. Les autres mesures de $(\vect{u},\vect{v})$ sont donc de la forme $x+2k\pi$. Or $\cos(x+2k\pi)=\cos x$ et $\sin(x+2k\pi)=\sin x$. On a donc la d\'efinition suivante :

\vspace{1ex}
\begin{defn}
Le cosinus (resp. le sinus) d'un angle orient\'e de vecteurs est le cosinus (resp. le sinus) d'une quelconque de ses mesures.
\end{defn}

\newpage%%%%%%%%%%%%%%%%%%%%%%%%%%%%%%%%%%%%%%%%%%%%%%%%%%%%%%%%%%%%%%%%%%%%%%%%%%%%%%%%%%%%%%%%%%%%%%%%%%%%
\subsection{Angles associ\'es}
L'utilisation de sym\'etries dans le cercle trigonom\'etrique permet d'\'etablir :% les relations suivantes :

\vspace{1ex}
\begin{prop}
Pour tout r\'eel $x$,

\vspace{1ex}
$\sysd{\cos(-x)=\cos x}{\sin(-x)=-\sin x}$ \hfill $\sysd{\cos(x+\pi)=-\cos x}{\sin(x+\pi)=-\sin x}$ \hfill $\sysd{\cos(\pi-x)=-\cos x}{\sin(\pi-x)=\sin x}$

\hfill $\sysd{\cos\left(\dfrac{\pi}{2}-x\right)=\sin x}{\sin\left(\dfrac{\pi}{2}-x\right)=\cos x}$ \hfill $\sysd{\cos\left(\dfrac{\pi}{2}+x\right)=-\sin x}{\sin\left(\dfrac{\pi}{2}+x\right)=\cos x}$ \hfill\phantom{.}

\includegraphics{CoursTrigoFigures.2} \hfill \includegraphics{CoursTrigoFigures.3}

\end{prop}



\subsection{Valeurs particuli\`eres}
à l'aide de consid\'erations g\'eom\'etriques, on peut obtenir les valeurs suivantes :

\renewcommand{\arraystretch}{3}
\begin{center}
\begin{tabular}{|c|c|c|c|c|c|c|} \hline
$x$ & 0 & $\dfrac{\pi}{6}$ & $\dfrac{\pi}{4}$ & $\dfrac{\pi}{3}$ & $\dfrac{\pi}{2}$ & $\pi$ \\ \hline
$\cos(x)$ & 1 & $\dfrac{\sqrt{3}}{2}$ & $\dfrac{\sqrt{2}}{2}$ & $\dfrac{1}{2}$ & 0 & -1 \\ \hline
$\sin(x)$ & 0 & $\dfrac{1}{2}$ & $\dfrac{\sqrt{2}}{2}$ & $\dfrac{\sqrt{3}}{2}$ & 1 & 0 \\ \hline
\end{tabular}
\end{center}
\renewcommand{\arraystretch}{1}
\subsection{Formules d'addition}
\begin{prop}
Quels que soient les nombres r\'eels $a$ et $b$,\\

\vspace{-1ex}
\begin{minipage}{0.48\linewidth}
\begin{equation}
\cos(a-b)=\cos a \cos b + \sin a \sin b
\end{equation}
\begin{equation}
\cos(a+b)=\cos a \cos b - \sin a \sin b
\end{equation}
\end{minipage}
\begin{minipage}{0.48\linewidth}
\begin{equation}
\sin(a-b)=\sin a \cos b - \cos a \sin b
\end{equation}
\begin{equation}
\sin(a+b)=\sin a \cos b + \cos a \sin b
\end{equation}
\end{minipage}
\end{prop}

\begin{demo}
\underline{\'equation (1.1)}: Soit $\oij$ un rep\`ere orthonormal, $\calc$ le cercle trigonom\'etrique de centre $O$ et $A$ et $B$ les points de $\calc$ tels que $(\vect{\imath},\vect{OA})=a$ et $(\vect{\imath},\vect{OB})=b$. 

\par Calculons de deux façons diff\'erentes le produit scalaire $\vect{OA}.\vect{OB}$\\

\begin{itemize}
\item $\vect{OA}.\vect{OB}=OA\times OB\times\cos(\vect{OA},\vect{OB})$
\par Or $(\vect{OA},\vect{OB})=(\vect{OA},\vect{\imath})+(\vect{\imath},\vect{OB})+2k\pi=-a+b+2k\pi$ et $OA=OB=1$
\par Ainsi $\vect{OA}.\vect{OB}=1\times1\times\cos(b-a)=\cos(a-b)$
\item Les coordonn\'ees de $\vect{OA}$ sont $\cp{\cos a}{\sin a}$ et les coordonn\'ees de $\vect{OB}$ sont $\cp{\cos b}{\sin b}$
\par On a donc : $\vect{OA}.\vect{OB}=\cos a \cos b + \sin a \sin b$
\end{itemize}

\underline{Conclusion }: $\cos(a-b)=\cos a \cos b + \sin a \sin b$\\

\underline{\'equation (1.2)}: $\cos(a+b)=\cos(a-(-b))=\cos a \cos (-b) + \sin a \sin (-b)=\cos a \cos b - \sin a \sin b$\\

\underline{\'equation (1.3) }: $\sin(a-b)=\cos\left( \dfrac{\pi}{2}-(a-b) \right)=\cos\left( \left(\dfrac{\pi}{2}-a \right)+b\right) = \cos\left(\dfrac{\pi}{2}-a \right)\cos b-\sin\left(\dfrac{\pi}{2}-a \right)\sin b = \sin a \cos b - \cos a \sin b$ \\

\underline{\'equation (1.4) }: $\sin(a+b)=\sin(a-(-b))=\sin a \cos (-b) - \cos a \sin (-b) = \sin a \cos b + \cos a \sin b$
\end{demo}


%%%%%%%%%%%%%%%%%%AJOUT AVRIL 2008%%%%%%%%%%%%%%%%%%%%%%%%%%%%%%%%%%%%%%%%%
\subsection{Formules de duplication}
\begin{prop}
Quel que soit le r\'eel $x$,
\begin{align}
 \cos(2x)&=\cos^2x-\sin^2x=2\cos^2x-1=1-2\sin^2x\label{cos2x} \\
 \sin(2x)&=2\sin x\cos x\label{sin2x}
\end{align}
\end{prop}

\begin{demo}
(\ref{cos2x}) : Quel que soit le r\'eel $x$,
 \[\cos(2x)=\cos(x+x)=\cos x\cos x-\sin x\sin x=\cos^2x-\sin^2x\]
En utilisant les relations $\sin^2x=1-\cos^2x$ et $\cos^2x=1-\sin^2x$, on obtient successivement :
\[\cos(2x)=2\cos^2x-1\quad\text{puis}\quad\cos(2x)=1-2\sin^2x\]
(\ref{sin2x}) : Quel que soit le r\'eel $x$,
 \[\sin(2x)=\sin(x+x)=\sin x\cos x+\cos x\sin x=2\sin x\cos x\]

\end{demo}

\section{R\'esolution d'\'equations}
\begin{prop}
Quels que soient les r\'eels $a$ et $b$,

%\begin{equation*}
 $\cos a=\cos b \Leftrightarrow
   \left|\text{ou }
     \begin{array}{@{}l}
	 a=b+k\times2\pi \\
	 a=-b+k\times2\pi
     \end{array}
   \right.$
     \hfill et\hfill
 $\sin a=\sin b \Leftrightarrow
   \left|\text{ou }
     \begin{array}{@{}l}
	 a=b+k\times2\pi \\
	 a=\pi-b+k\times2\pi
     \end{array}
   \right.$
%\end{equation*}
\end{prop}

\begin{exemple}
R\'esoudre dans $]-\pi~;~\pi[$ l'\'equation $\sin\left(x-\frac{\pi}{3}\right)=\frac{1}{2}$
\end{exemple}
\begin{align*}
\sin\left(x+\frac{\pi}{3}\right)=\frac{1}{2} &\Leftrightarrow \sin\left(x-\frac{\pi}{3}\right)=\sin \frac{\pi}{6}\\
&\Leftrightarrow
    \left|\text{ou }
     \begin{array}{@{}l}
	 x-\frac{\pi}{3}=\frac{\pi}{6}+k\times2\pi \\
	 x-\frac{\pi}{3}=\pi-\frac{\pi}{6}+k\times2\pi
     \end{array}
   \right.\\
&\Leftrightarrow
    \left|\text{ou }
     \begin{array}{@{}l}
	 x=\frac{\pi}{2}+k\times2\pi \\
	 x=\frac{7\pi}{6}+k\times2\pi
     \end{array}
   \right.
\end{align*}
Sur l'intervalle $]-\pi~;~\pi[$, on a $\cals=\{\frac{\pi}{2};-\frac{5\pi}{6}\}$

%%%%%%%%%%%%%%%%%%FIN AJOUT%%%%%%%%%%%%%%%%%%%%%%%%%%%%%%%%%%%%%%%%%%%%%%%%%

\newpage
\section{Rep\'erage polaire}
\subsection{Coordonn\'ees polaires d'un point}
\begin{minipage}{11cm}
Soit $\oij$ un rep\`ere orthonormal direct. Si $M$ est un point distinct du point $O$ alors $M$ peut-\^etre rep\'er\'e par l'angle $\theta=(\vect{\imath},\vect{OM})$ et la longueur $r=OM$.
\par R\'eciproquement, la donn\'ee d'un couple $(r;\theta)$ avec $r>0$ et $\theta\in\R$ d\'etermine un seul point $M$ tel que $\theta=(\vect{\imath},\vect{OM})$ et $r=OM$.
\end{minipage}
\hfill
\begin{minipage}{3.2cm}
\includegraphics[scale=1.3]{CoursTrigoFigures.4}
\end{minipage}

%\vspace{1ex}
\begin{defn}
Pour tout point $M$ distinct de $O$, un couple $(r;\theta)$ tel que $\theta=(\vect{\imath},\vect{OM})$ et $r=OM$ est appel\'e couple de coordonn\'ees polaires de $M$ dans le rep\`ere polaire $(O,\vect{\imath})$.
\end{defn}


\subsection{Lien entre coordonn\'ees cart\'esiennes et polaires}
Soit $\oij$ un rep\`ere orthonormal direct.

\begin{prop}
Si $M$ est un point ayant pour coordonn\'ees cart\'esiennes $(x;y)$ dans le rep\`ere $\oij$ et pour coordonn\'ees polaires $(r;\theta)$ alors :
\begin{center}
$r=\sqrt{x^{2}+y^{2}}$ \qquad  ; \qquad $x=r\cos \theta$ \qquad ; \qquad $y=r\sin \theta$
\end{center}
\end{prop}

\begin{demo}
Notons $\calc$ le cercle trigonom\'etrique de centre $O$. La demi-droite $[OM)$ coupe $\calc$ en $N$. On a donc $\vect{OM}=r\vect{ON}$.
\par $N\in\calc$ donc ses coordonn\'ees cart\'esiennes sont $(\cos \theta ; \sin \theta)$. Celles de $M$ sont donc $(r\cos \theta ; r\sin \theta)$.
\par Par unicit\'e des coordonnées, $x=r\cos \theta$ et $y=r\sin \theta$.

De plus $OM^{2}=x^{2}+y^{2}$ et $OM=r$ donc $r^{2}=x^{2}+y^{2}$.
\end{demo}





\end{document}
